% !TeX root = ../教材题库_全部条目_转写工作稿.tex
% Source: XJ-B1-C06-EX-P244-Y0613
% 第6章 / 6.4 用样本估计总体 / 教材 p.244、245、246 / PDF 第251、252、253页
% OCR未校对：公式与文字均可能有误；未完成数学TeX转写；禁止直接发给学生
% 原题及相邻原解截图（核对用）：../题目原图/XJ-B1-C06-EX-P244-Y0613-01.webp, ../题目原图/XJ-B1-C06-EX-P244-Y0613-02.webp, ../题目原图/XJ-B1-C06-EX-P244-Y0613-03.webp
% 本文件只在转写完成后作为题干使用；OCR 条目尚不是可靠题干。
\begin{bankitem}{6.4 例题 例题7}{第6章 / 6.4 用样本估计总体 / 教材 p.244、245、246 / PDF 第251、252、253页}{XJ-B1-C06-EX-P244-Y0613}{OCR未校对：公式与文字均可能有误；未完成数学TeX转写；禁止直接发给学生}
\begin{ocrdraft}
[教材p.244 / PDF第251页]
\par
0才树人中学高一年级有男生180人，女生120人．该校高一年级统计小组
\par
为调查本年级学生身高情况，采取分层抽样的方法从男生和女生中分别抽取18人，
\par
12人．将男生组看作样本A，计算出样本AI的平均身高为173．5cm,方差为17；
\par
将女生组看作样本A2，计算出样本A2的平均身高为164．0cm，方差为30．试根据
\par
以上数据计算由A,A2组成的样本A的方差，并估计总体方差．
\par
分析按分层抽样获取的样本分为两层：男生组与女生组．现已知男生组样本
\par
和女生组样本的均值与方差，借助方差的定义可计算出分层抽样样本的方差，进而
\par
估计总体方差．
\par
解设从男生中抽出的样本个体为YI，，\ensuremath{\cdot}\ensuremath{\cdot}\ensuremath{\cdot}，Y18，均值记为3，，方差记为《；
\par
从女生中抽取的样本个体为之
\par
均值记为之，方差记为
\par
\ensuremath{\cdot}，之12，
\par
先计算样本A的均值
\par
十十'十308十0十之2十\ensuremath{\cdot}
\par
18十12
\par
18T十12之
\par
18十12
\par
18\ensuremath{\times}173．5十12\ensuremath{\times}164．0
\par
30
\par
=169．7(cm)；
\par
再计算样本A的方差：
\par
1
\par
一[（一）2十()2一2十．'十(Y18一）2十
\par
30
\par
12
\par
1
\par
30
\par
18
\par
1
\par
一到十
\par
30
\par
[教材p.245 / PDF第252页]
\par
2十一的
\par
1
\par
18
\par
30
\par
2十一的
\par
12
\par
12
\par
1
\par
30
\par
1
\par
一08《十18(y一）2十12十12（一
\par
30
\par
1
\par
．一（18\ensuremath{\times}17十18\ensuremath{\times}3．82十12\ensuremath{\times}30十12\ensuremath{\times}5，72）
\par
30
\par
=43．86\ensuremath{\cdot}
\par
于是可以估计该校高一年级学生身高的方差为43，86\ensuremath{\cdot}
\par
分层抽样在获得总体方差估计的同时，也得到各层的方差估计．
\par
从例7可知，如果将总体分为两层，第一、二层的样本量分别为，\ocrglyph{2116}，样本
\par
均值分别为云，。，样本方差分别为则全部样本的样本容量、样本均值和
\par
样本方差分别为
\par
刀==771H\ocrglyph{00AF}刀2，
\par
\ocrglyph{00AF}(n1X1\ocrglyph{00AF}\ocrglyph{00E5}一772T2），
\par
如果将总体分为差e
\par
Z，差冫3）层，你能猜测出
\par
一\{时十0一2]十十02
\par
1
\par
结果吗？试一试．
\par
0
\par
标准差
\par
方差充分利用所有数据，并且仅用一个数值来刻画一组数据的离散程度，但方
\par
差也有局限性，如方差的单位是观测数据的单位的平方，而刻画离散程度的一种理
\par
想度量应当具有与观测数据相同的单位．解决这一局限性的方法就是引人标准差．
\par
标准差是方差的算术平方根．
\par
如果乎是总体方差，则称。一、，是总体标准差；
\par
如果0是样本方差，则称s一，/0是样本标准差．
\par
给定数据
\par
\ensuremath{\cdot}\ensuremath{\cdot}，葑和均值由方差计算公式知道，样本标准差5可以用
\par
下面的公式计算：
\par
1
\par
[教材p.246 / PDF第253页]
\par
显然，在刻画观测数据的离散程度上，方差与标准差是一样的．在对实际问题
\par
进行分析时，人们更多地使用标准差．
\par
样本标准差依赖于样本的选取，也带有随机性．样本方差是总体方差的估计．
\par
样本标准差是总体标准差的估计．
\end{ocrdraft}
\end{bankitem}
